\documentclass[12pt]{article}
\usepackage{fullpage}
\usepackage{times}
\usepackage{fancyhdr,graphicx,amsmath,amssymb}
\usepackage[ruled,vlined]{algorithm2e}
\usepackage[active,tightpage]{preview}

\PreviewEnvironment{algorithm}
\begin{document}

\begin{algorithm}[H]
\SetAlgoLined
\KwIn{Dataset \( X \in \mathbb{R}^{n \times M} \), neighborhood radii \( \{ \epsilon_1, \epsilon_2, \ldots, \epsilon_N \} \), minimum number of samples \texttt{minPts} to form a dense region.}
\KwOut{Outlierness scores vector \(S \in \mathbb{R}^n \) with normalized outlierness values. Each value is within \([0; 1]\) range. Values closer to 1 indicate stronger outliers.}

Initialize outlierness scores vector \(S \leftarrow \mathbf{0} \in \mathbb{R}^n\)\;

Build spatial index (k-d tree) on dataset \( X \)\;

\For{\( i \leftarrow 1 \) \KwTo \( N \)}{
    Given parameters \( \epsilon = \epsilon_i \) and \texttt{minPts}:
    
    % \tcp{Simplified DBSCAN steps:}
    \Indp
    Identify core points with at least \texttt{minPts} neighbors within radius \(\epsilon_i\)\;
    
    Expand clusters from core points (without assigning explicit cluster labels)\;

    Assign binary outlier labels \( l_i \in \{0,1\}^n \), where 1 denotes an outlier\;
    \Indm
    
    Update scores: \(S \leftarrow S + l_i\)\;
}

Normalize scores by number of radius values:  
\[
S \leftarrow \frac{S}{N}
\]

\Return \(S\);

\caption{DBSOD: Density-Based Spatial Outlier Detection}
\end{algorithm}

\end{document}
